\documentclass[10pt,a4paper]{amsart}
\usepackage[margin=19mm]{geometry}
\usepackage{amsmath,amssymb,mathtools}
\usepackage[scr=rsfs,cal=euler]{mathalfa}
\usepackage{xfrac}
\usepackage[unicode,hidelinks,bookmarks=false]{hyperref}
\newcommand{\ie}{i.e.\ }
\newcommand{\cf}{cf.\ }
\newcommand{\resp}{respectively\ }
\newcommand{\Subsection}{\subsection}
\providecommand{\nobreakdash}{\nobreak\hbox{-}}
\setlength{\emergencystretch}{2em}
\multlinegap0pt
\usepackage{bm}
\usepackage{bbm}
\usepackage{dsfont}
\usepackage{xspace}
\usepackage{cancel}
\usepackage{mathtools}
\usepackage{amsthm}
\makeatletter
\@ifundefined{lem}{\newtheorem{lem}{Lem}}{}
\@ifundefined{proposition}{\newtheorem{proposition}[lem]{Proposition}}{}
\makeatother
\newcommand{\Z}{\mathbb{Z}}
\title[A classification algorithm for reflexive simplices]{A classification algorithm for reflexive simplices}
\author{Marco Ghirlanda}
\date{Source: arXiv:2510.09131v1 (CC BY 4.0)}
\begin{document}
\maketitle

\noindent\emph{Source and licence.} A classification algorithm for reflexive simplices. Version arXiv:2510.09131v1 (CC BY 4.0), distributed under the Creative Commons Attribution 4.0 International licence, as recorded in the arXiv licence metadata for this exact version and shown on the abstract page https://arxiv.org/abs/2510.09131v1. Full rights notice: licenses/arxiv-2510.09131-CC-BY-4.0.txt.\par\smallskip

\noindent\textit{Editorial scope.} Counted unit: the Proposition whose source label is \texttt{generation condition} in the article (source0.tex lines 306--317, source label \texttt{generation condition}), delivered together with the article's own complete original proof of it (source0.tex lines 319--360, one \texttt{proof} environment of 42 lines). No mathematics is written, repaired or re-derived: statement and proof are the article's own text line for line. Required context retained uncounted: none. Every definition, display and label of those lines is retained unchanged. Omitted from this package and not redistributed: 1--305, 318--318, 361--960. Nothing inside a retained excerpt is omitted or altered: every retained body is delivered exactly as the article's lines are written, so no omission receipt of any kind accompanies this package. Citations: the retained excerpts carry no citation command at all, so no bibliography entry of the article is transcribed and no reference is redistributed. Reference adaptation: none was needed and none was made; every reference inside the delivered unit resolves inside it, so no unresolved reference or citation is produced. Printed slips kept literal: no printed word, symbol, hypothesis or number of the counted statement or of its proof was altered. Layout and class: the article's own class is \texttt{amsart} with packages including inputenc, babel, amsmath, amsfonts, amssymb, amsthm, listings, hyperref, mathtools, xy, youngtab, array, csquotes, tikz-cd; the wrapper is the lane's pinned \texttt{amsart} editorial wrapper at 10pt on A4, which restates the theorem-like environments the retained text needs and adds the article's own macro definitions that the retained text needs (\texttt{\textbackslash Z}), with extra wrapper packages (bm, bbm, dsfont, xspace, cancel, mathtools). The delivered numbering is the unit's own. Assets: no retained span contains a figure environment, a graphics-inclusion command, a PDF-page inclusion, a table or a TikZ picture; no image file is copied into this package, the package declares no asset dependency and no image file is redistributed with it. Licence: version arXiv:2510.09131v1 of this article is distributed under the Creative Commons Attribution 4.0 International licence, as recorded in the arXiv licence metadata for this exact version and shown on the abstract page https://arxiv.org/abs/2510.09131v1. A full rights notice is delivered at licenses/arxiv-2510.09131-CC-BY-4.0.txt. Characters: every retained excerpt is pure ASCII, so the delivered engine, XeLaTeX with its default Latin Modern text font, reports no missing character.
\begin{proposition}\label{generation condition}
    Let $G=\Z^k\oplus \Z/\mu_1\Z\oplus\dots\oplus\Z/\mu_r\Z$ be in invariant factor form, and let $\omega_i=(w_i,\eta_i)\in G$ for $i=1,\dots,n$. For $j=0,\dots,r$ set
    $$
    Q_j\coloneqq \begin{bmatrix}
        w_1 &\dots &w_n\\
        \eta_{11} & \dots & \eta_{n1}\\
        \vdots & & \vdots\\
        \eta_{1j} & \dots & \eta_{nj}
    \end{bmatrix}.
    $$
    Then $G=\langle \omega_1,\dots,\omega_n\rangle$ if and only if the maximal minors of $Q_0$ generate $\Z$ and the maximal minors of $Q_j$ generate $\Z/\mu_j\Z$ for any $j=1,\dots,r$. 
\end{proposition}

\begin{proof}
    Let $G=\langle \omega_1,\dots,\omega_n\rangle$. By classical determinantal divisors theory, the maximal minors of $Q_0$ generate $\Z$. Fix $j=1,\dots,r$, and consider the prime factorization~$\mu_j=p_1^{m_1}\cdots p_t^{m_t}$. Now, fix $i=1,\dots,t$, and consider the projection onto the first $k+j$ coordinates followed by the coordinate-wise projection onto $\Z/p_i\Z$:
    $$
    \pi\colon G\rightarrow(\Z/p_i\Z)^{k+i}.
    $$
    Then $\pi(\omega_1),\dots,\pi(\omega_n)$ forms a basis of the vector space $(\Z/p_i\Z)^{k+j}$, encoded by the matrix $Q_j$ modulo $p_i$. In particular, there exists a maximal minor of $Q_j$ generating~$\Z/p_1\Z$, and hence $\Z/p_1^{m_1}\Z$. The assertion follows since, by the Chinese remainder theorem, we have an isomorphism
    $$
    \Z/\mu_j\Z\rightarrow\Z/{p_1^{m_1}}\Z\oplus\dots\oplus\Z/p_t^{m_t}\Z, \quad     \eta\mapsto(\eta\mod p_1^{m_1},\dots,\eta\mod p_t^{m_t}).
    $$
    
    Now assume that the maximal minors of $Q_0$ generate $\Z$ and the maximal minors of $Q_j$ generate $\Z/\mu_j\Z$ for any $j=1,\dots,r$.
    Let $\eta'_{ij}\in\Z_{\geq 0}$ be the representing integer of  $\eta_{ij}$ between $0$ and $\mu_j-1$, and consider the integer matrix
    $$
    Q'_j\coloneqq \begin{bmatrix}
        w_1 &\dots &w_n\\
        \eta'_{11} & \dots & \eta'_{n1}\\
        \vdots & & \vdots\\
        \eta'_{1j} & \dots & \eta'_{nj}
    \end{bmatrix}.
    $$
    Denote by $q_j$ the greatest common divisor of all maximal minors of $Q'_j$. By assumption, we have $q_0=1$ and $\gcd(q_j,\mu_j)=1$.
    With the vectors~$b_i\coloneqq\mu_i e_{k+i}$ we define the $(k+r)\times(n+r)$-matrix
    $$
    M\coloneqq[Q'_r,B],\qquad B\coloneqq[b_1,\dots,b_r].
    $$
    Denote by $\Delta$ the greatest common divisor of all maximal minors of $M$. Notice that~$\Delta$ divides $\mu_r\cdots\mu_{i+1}q_{i}$ for all $i=0,\dots,r$.
    In particular, $\Delta$ divides \begin{align*}
        &\gcd(q_r,\mu_r q_{r-1},\mu_r\mu_{r-1}q_{r-2},\dots,\mu_r\cdots\mu_1q_0)\\
        =&\gcd(q_r,\mu_r\gcd(q_{r-1},\mu_{r-1}q_{r-2},\dots,\mu_{r-1}\cdots\mu_1q_0)).
    \end{align*}
    Since $\gcd(q_r,\mu_r)=1$, we have
    \begin{align*}
    &\gcd(q_r,\mu_r\gcd(q_{r-1},\mu_{r-1}q_{r-2},\dots,\mu_{r-1}\cdots\mu_1q_0))\\
    =&\gcd(q_r,q_{r-1},\mu_{r-1}q_{r-2},\dots,\mu_{r-1}\cdots\mu_1q_0).
    \end{align*}
    Iterating the argument for $\mu_{r-1},\dots,\mu_1$, we see that $\Delta$ divides~$\gcd(q_r,\dots,q_0)$. Since~$q_0=1$, it follows that $\Delta=1$. By classical determinantal divisors theory, it follows that
    $\Z^{k+r}=\langle M_{*,1},\dots,M_{*,n+r}\rangle$. Thus, the equality~$G=\langle \omega_1,\dots,\omega_n\rangle$ follows from
    $$
    \dfrac{G}{\langle \omega_1,\dots,\omega_r\rangle}\cong\dfrac{\Z^{k+r}}{\langle M_{*,1},\dots,M_{*,n+r}\rangle}.
    $$
    
\end{proof}

\end{document}
